\defmodule{coeff}

\baselineskip=14pt
Ce programme permet de trouver des bons coefficients pour un
g\'en\'erateur dont on a donn\'e le modulo et le type. Cette
recherche se fait en fonction des $P_\alpha$. Il faut donc 
\'egalement pr\'eciser un $\alpha$. 

\noindent
{\bf Le format du fichier d'entr\'ee pour Coeff}

\begin{center}
\it{
\begin{tabbing}
\hspace{2cm}\= LCG \hspace{3cm} \= SorteGenerateur \\
\>Coeff \>   Type de calcul \\
\>251     \>      m  modulo\\
\>true    \>      modulo premier\\  
\>3 250   \>      intervalle pour le multiplicateur a\\
\>2       \>      mindim\\
\>4      \>      maxdim\\
\>2      \>      saut \\ 
\>2       \>      alpha\\
\>false   \>      VerifBooleen\\
\>1       \>      seed\\
\>1.0 2.0 3.0 1.0 1.0 \> Beta(0),...,Beta(MaxDim)\\
\end{tabbing}
}
\rm {Figure 1 : Format du fichier pour Coeff}
\end{center}

\noindent
Tous les champs sont identiques au cas du PAl ou BAl sauf pour le
cinqui\`eme qui indique plut\^ot l'intervalle \`a l'int\'erieur 
duquel chercher le meilleur coefficient.
